\documentclass[11pt]{article}
\usepackage[T1]{fontenc}
\usepackage[utf8]{inputenc}
\usepackage{lmodern,amsmath,amssymb,amsthm,booktabs,array}
\usepackage[margin=1in]{geometry}
\usepackage[colorlinks=true,urlcolor=blue,citecolor=blue,linkcolor=blue]{hyperref}
\usepackage{microtype}
\newtheorem{theorem}{Theorem}
\newtheorem{lemma}[theorem]{Lemma}
\newtheorem{corollary}[theorem]{Corollary}
\newcommand{\Z}{\mathbb Z}
\newcommand{\N}{\mathbb N}
\newcommand{\leanfile}[2]{\href{https://github.com/the-omega-institute/trureturing/blob/f960cb09514c99f42fd07398f5240a77378ea4db/D5/S3/Combinatorics/Graph/Roommates/#1.lean}{#2}}
\title{Mertens's range-two stable-matching conjecture\protect\\and fully supported preference laws}
\author{Haobo Ma\\\small ChronoAI Pte Ltd, Singapore
\and Wenlin Zhang\\\small National University of Singapore\\\small ChronoAI Pte Ltd, Singapore}
\date{November 16, 2026}
\begin{document}
\maketitle
\begin{abstract}
For Mertens's Conjecture~3 in \emph{Random stable matchings}
\cite[Section~3.2, equation~(15), p.~7; arXiv:cond-mat/0509221v2]{mertens2005},
we prove $p_N=\Theta(q^N)$ for the one-dimensional range-two model,
simultaneously for both parities, by a six-site completion of open
preference words. Here $p_N$ is the probability of a stable partial
matching, and $q\in(0,1)$. The same construction proves this conclusion
for every fully supported independent identically distributed law on the
24 local rankings. For the uniform law it gives explicit finite-size
bounds and convergent enclosures for $q$. We also give the elementary
range-one comparison and triangle-packing bounds, and explain the
complete-graph endpoint of Conjecture~5 using Young and Studebaker's
corrected asymptotic constant.
\end{abstract}

\section{The probability and its exponential base}
Mertens predicted exponential decay of the solvability probability on
periodic finite-range grids \cite[Conjecture~3]{mertens2005}.
For range two in one dimension, the local exclusive-triangle obstruction
already gives exponential upper bounds, but it does not identify a base
that also gives a lower bound up to a constant factor. We obtain that
common base by a six-site completion of open intervals.

For $N\geq5$, take the labelled vertex set $\Z/N\Z$ and join each vertex
$i$ to $i+1,i+2,i-1,i-2$. These four neighbors are distinct. Independently
at each vertex choose a uniform strict ranking of its four neighbors;
there are $24^N$ preference profiles. Every neighbor is preferable to
being single. A partial matching is stable when no edge outside the
matching has both endpoints preferring one another to their assignments.
Thus a single endpoint always prefers an acceptable partner, and two
adjacent singles block. This is the convention of
\cite[Section~2]{mertens2005}, not a perfect-matching requirement.
Write $A_N$ for the number of profiles with at least one stable partial
matching and $p_N=A_N/24^N$. Each profile is counted once, regardless of
its number of stable matchings.

For the auxiliary open model let $I_m=\{0,\ldots,m-1\}$, with $I_0$
empty. A ranking word of length $m$ assigns to each $i\in I_m$ a
bijection $r_i:\{-2,-1,1,2\}\to\{0,1,2,3\}$, smaller ranks being
preferred. Set $r_i(0)=4$. A certificate is a displacement vector
$d:I_m\to\{-2,-1,0,1,2\}$ with the following properties:
\begin{enumerate}
\item If $d_i\ne0$ and $i+d_i=j\in I_m$, then $d_j=-d_i$.
\item The targets $i+d_i$ with $d_i\ne0$ are distinct integers,
including targets outside $I_m$.
\item For $i<j\leq i+2$ in $I_m$, either $d_i=j-i$, or
\[
 r_i(d_i)<r_i(j-i)\quad\hbox{or}\quad r_j(d_j)<r_j(i-j).
\]
\end{enumerate}
An exterior request retains its actual neighbor rank; it is not treated
as being single. Let $b_m$ be the number of words admitting at least one
certificate. Thus $b_0=1$ and $b_m\leq24^m$.

\begin{theorem}\label{thm:main}
There is a number $\lambda$ satisfying
\[
  6\leq\lambda\leq 13808^{1/3}<24
\]
such that, on setting $q=\lambda/24$, every integer $N\geq7$ satisfies
\begin{equation}\label{eq:main}
 \lambda^{-6}q^N\ \leq\ p_N\ \leq\ b_6q^N,
 \qquad 1\leq b_6\leq24^6.
\end{equation}
In particular,
\[
 \frac14\leq q\leq\left(1-\frac1{864}\right)^{1/3}<1,
 \qquad \lim_{N\to\infty}p_N^{1/N}=q.
\]
The same base works for both parities.
\end{theorem}

The estimate is qualitative: it does not determine $q$ or establish the
numerical fit near $0.97$ in \cite[Section~3.2]{mertens2005}.
The complete-graph problem has a different neighborhood and asymptotic
regime. Mertens's later simulations \cite{mertens2015} address that regime;
Chin and Michelen \cite[Theorem~1]{chin2026} give a vanishing
upper bound in that complete-graph model; subsequently Young and
Studebaker \cite[Theorem~1.1]{young2026} prove its $N^{-1/4}$
asymptotic with a corrected leading constant. Neither statement replaces
the fixed-degree cyclic estimate proved here.

\section{Restriction and a six-site completion}

\begin{lemma}\label{lem:restriction}
For $m,k\geq0$,
\[
 6^m\leq b_m,\qquad b_{m+k}\leq b_mb_k.
\]
\end{lemma}
\begin{proof}
Restriction to an interval preserves all three certificate conditions:
translate both sites and targets by the interval's starting point.
A partner removed by restriction becomes an exterior request with its
original displacement and rank. Splitting a word of length $m+k$ into
its two ranking subwords is injective, so $b_{m+k}\leq b_mb_k$.

For the lower bound prescribe displacements $+1,-1,+1,-1,\ldots$.
When $m$ is odd the final $+1$ is an exterior request. Targets are
distinct and all internal requests are reciprocal. Put the prescribed
displacement first at every site. All sites then have their first
choice, so no internal edge blocks. Each site has $3!=6$ ways to rank
the remaining offsets, independently.
\end{proof}

\begin{lemma}[Six-site completion]\label{lem:completion}
For $m,k\geq1$,
\begin{equation}\label{eq:injections}
 b_mb_k\leq b_{m+k+6},\qquad b_m\leq A_{m+6}.
\end{equation}
Also $A_N\leq b_N$ for every $N\geq5$.
\end{lemma}
\begin{proof}
First place two certified words on $[0,m-1]$ and
$[m+6,m+6+k-1]$, leaving six sites between them. Label these sites
$0,1,2,3,4,5$ locally. Requests from the left word into the gap land at
sites $0,1$; requests from the right word land at $4,5$. Injectivity of
targets makes the requests distinct within each word, and the two
landing sets are disjoint. Fulfill every such request by matching its
old endpoint to its requested gap site.

The remaining free sites of $\{0,1,2\}$ form a clique containing $2$;
the remaining free sites of $\{3,4,5\}$ form a clique containing $3$.
In each clique of even size pair all its sites. In a clique of odd size,
leave its center single and pair the other sites. Every such pair has
distance at most two. If both centers remain single, match $2$ to $3$.
Consequently at most one gap site remains single, and it is $2$ or $3$.

At each matched gap site put its assigned partner first and order its
other offsets arbitrarily. Rank offsets at a single gap site arbitrarily.
Old displacements, including requests at the two outer ends of the
whole interval, remain unchanged. Reciprocity and target injectivity
hold: an old site requesting the gap cannot already be the target of
another old site, by its original reciprocal condition; reserved gap
sites are distinct and are excluded from the free-site pairing.
The two old intervals are more than distance two apart, so their targets
cannot collide and they have no edge between them.

There is no new blocking edge. Edges inside an old word retain the
original inequalities. A matched gap endpoint has its first choice and
rejects every alternative. A single gap site has no old neighbor: its
distance from either old interval is at least three. Every gap neighbor
of that single is matched, since there is at most one single in the
gap. These observations cover every edge of the completed word.

For definiteness choose the lexicographically least certificate for each
input word and fixed choices for the finite gap construction. The map
on words is injective because restriction to the two fixed old intervals
recovers both inputs. No matching witnesses are counted. This proves
the first inequality in \eqref{eq:injections}.

For the second inequality place one word on $[0,m-1]$ in a cycle of
length $m+6$ and append six gap sites. Its right exterior targets
$m,m+1$ occupy gap sites $0,1$; its left targets $-2,-1$ become gap
sites $4,5$ after reduction modulo $m+6$. These four targets are
distinct; only the targets actually requested are occupied. The same
completion applies. The old ranks of signed neighbors are unchanged.
No edge connects old sites across the six-site gap, and gap centers
remain at cyclic distance at least three from every old site. Thus the
same edge check proves stability, even when $m=1$ or $m=2$.
Restriction recovers the original ranking word, proving injectivity.

Finally cut a stable cyclic matching at the label $0$. Assign displacement
$0$ to a single, and the unique signed displacement in $\{-2,-1,1,2\}$
to every matched vertex. Uniqueness holds already at $N=5$.
Reciprocity survives for partners inside the interval. Equality of two
integer targets would imply equality of their cyclic targets, contradicting
injectivity of the matching. Every internal edge retains its original
nonblocking condition. The resulting word is certified; retaining the
whole ranking profile gives an injection and $A_N\leq b_N$.
\end{proof}

\section{Strict loss and the common growth rate}
\begin{lemma}\label{lem:triangle}
$b_3\leq24^3-16=13808$.
\end{lemma}
\begin{proof}
On three consecutive sites put the two internal neighbors first and
second in cyclic order: each vertex ranks its predecessor first and its
successor second. If a certificate contains no internal pair, every
vertex prefers either internal neighbor to its exterior request or to
being single, so an internal edge blocks. If there is an internal pair,
the remaining vertex prefers both members of that pair to its assignment.
One member of the pair ranks this remaining vertex first and its
assigned partner second, so their edge blocks. Thus no certificate exists.
This is the exclusive odd-cycle obstruction of
\cite[Section~2 and equation~(14)]{mertens2005}.

There are two cyclic orientations. Once one is fixed, the two exterior
offsets can be ordered in $2!$ ways at each of three sites, giving
$2\cdot2^3=16$ distinct forbidden words. Their orientations differ in
their first-ranked internal neighbors, so the two sets are disjoint.
\end{proof}

\begin{lemma}\label{lem:defect}
There is $h\in[\log6,\log24)$ such that, for every $m\geq0$,
\begin{equation}\label{eq:defect}
 hm\leq\log b_m\leq hm+\log b_6.
\end{equation}
Moreover $3h\leq\log13808$.
\end{lemma}
\begin{proof}
Put $f_m=\log b_m$ and $D=f_6$. By Lemmas~\ref{lem:restriction}
and~\ref{lem:completion},
\[
 f_{m+k}\leq f_m+f_k,\qquad
 f_m+f_k\leq f_{m+k+6}\leq f_{m+k}+D \quad(m,k\geq1).
\]
The reverse inequality also holds when one index is zero because
$f_0=0$ and $D\geq0$. Define $h=\inf_{j\geq1}f_j/j$.
The lower bound $b_j\geq6^j$ gives $h\geq\log6$.
For fixed $j$, divide $m=aj+r$, $0\leq r<j$, to obtain
$f_m\leq af_j+f_r$. The finitely many remainders have bounded $f_r$,
so $\limsup f_m/m\leq f_j/j$. Taking the infimum and using the
definition of $h$ gives $f_m/m\to h$ and $hm\leq f_m$.

For fixed $m\geq1$, iteration of the reverse inequality gives
$f_{am}\geq af_m-(a-1)D$. Divide by $am$ and let $a\to\infty$:
$h\geq(f_m-D)/m$. This is the upper bound in \eqref{eq:defect};
the case $m=0$ follows from $D\geq0$.
Finally $3h\leq f_3\leq\log13808<3\log24$.
\end{proof}

\begin{proof}[Proof of Theorem~\ref{thm:main}]
Let $\lambda=e^h$. Exponentiating \eqref{eq:defect} gives
$\lambda^m\leq b_m\leq b_6\lambda^m$. For $N\geq7$,
Lemma~\ref{lem:completion} supplies
\[
 \lambda^{N-6}\leq b_{N-6}\leq A_N\leq b_N\leq b_6\lambda^N.
\]
Divide by $24^N$ to get \eqref{eq:main} with $q=\lambda/24$.
The bounds on $h$ give the bounds on $\lambda$ and $q$.
Taking $N$th roots in \eqref{eq:main} proves the limit along all
integers, not just one parity.
\end{proof}

\begin{corollary}\label{cor:prefix}
The same $q$ admits positive constants $c_*,C_*$ such that
$c_*q^N\leq p_N\leq C_*q^N$ for every $N\geq5$.
\end{corollary}
\begin{proof}
At $N=5$ match $0$ to $1$ and $2$ to $3$, leaving $4$ single.
Put the assigned partner first at each matched vertex. Any ranking at
the single vertex then works, since all its neighbors are matched to
their first choice. These profiles have probability $4^{-4}$.
At $N=6$ prescribe the adjacent pairs $(0,1),(2,3),(4,5)$ and put
each partner first, giving probability $4^{-6}$.
Thus $p_5\geq4^{-4}$ and $p_6\geq4^{-6}$. One may take
\[
 c_*=\min\{\lambda^{-6},4^{-4}q^{-5},4^{-6}q^{-6}\},\qquad
 C_*=\max\{b_6,q^{-5},q^{-6}\},
\]
using $p_5,p_6\leq1$ for the two upper bounds.
\end{proof}

\section{Consequences and neighboring source statements}
At $N=5$ the range-two cycle is the complete graph $K_5$, and every
stable partial matching has exactly one single: the number of singles is
odd, and any two singles would block. Thus Mertens's exact computation
\cite[Section~5, equation~(21b)]{mertenssmall2015} gives
$p_5=4075/6912$, strengthening the elementary bound $p_5\geq4^{-4}$
in Corollary~\ref{cor:prefix}.

The first consequence makes the common base accessible through finite
open-word counts, without a parity restriction.
\begin{corollary}\label{cor:enclosures}
For every $m\geq1$,
\[
 \frac{(b_m/b_6)^{1/m}}{24}\leq q\leq\frac{b_m^{1/m}}{24}.
\]
The logarithms of the endpoints differ by at most $6\log(24)/m$.
In particular $q$ is a computable real number, and both
$p_{2j}^{1/(2j)}$ and $p_{2j+1}^{1/(2j+1)}$ converge to $q$.
\end{corollary}
\begin{proof}
Rearrange $\lambda^m\leq b_m\leq b_6\lambda^m$ and use
$b_6\leq24^6$. For fixed $m$, the certificate definition is a finite
decision procedure, so $b_m$ and $b_6$ are exact computable integers.
The endpoint ratio tends effectively to one. Both subsequence limits
follow from Theorem~\ref{thm:main}.
\end{proof}

\begin{corollary}[Fully supported rankings]\label{cor:weighted}
Give each of the 24 rankings $\rho$ a probability $w(\rho)>0$, the same
at every site, with independent choices. Put
$\alpha=\min_\rho w(\rho)$, and let $a$ and $b$ be the probabilities
of ranking offset $+1$ and offset $-1$ first, respectively.
Let $\delta$ be the probability that three consecutive sites form either
orientation of an exclusive triangle. There is a common base $q_w$ with
\[
 \sqrt{ab}\leq q_w\leq(1-\delta)^{1/3}<1
\]
such that the solvability probability $p_N(w)$ satisfies, for $N\geq7$,
\[
 \alpha^6 q_w^{-6}q_w^N\leq p_N(w)\leq\alpha^{-6}q_w^N.
\]
Positive constants extend the same conclusion to every $N\geq5$.
This generalization is not posed in the source.
\end{corollary}
\begin{proof}
Let $B_m$ be the total product probability of certified open words,
with $B_0=1$. The restriction and six-site maps give
\[
 B_{m+k}\leq B_mB_k,\qquad
 \alpha^6B_mB_k\leq B_{m+k+6},\qquad
 \alpha^6B_{N-6}\leq p_N(w)\leq B_N.
\]
Indeed each selected completion contributes the original product weight
times six factors at least $\alpha$, and the maps are injective.
Since $B_{m+k+6}\leq B_{m+k}B_6\leq B_{m+k}$, the logarithms
$g_m=\log B_m$ obey
$g_m+g_k-6\log(1/\alpha)\leq g_{m+k}\leq g_m+g_k$.
The proof of Lemma~\ref{lem:defect} applies without the condition
$g_m\geq0$: it uses only a finite linear lower bound, supplied here by
$B_m\geq a^{\lceil m/2\rceil}b^{\lfloor m/2\rfloor}$.
Thus for $h_w=\inf_{m\geq1}g_m/m$,
\[
 e^{h_wm}\leq B_m\leq\alpha^{-6}e^{h_wm}.
\]
The displayed lower bound implies $e^{h_w}\geq\sqrt{ab}$.
The exclusive triangles exclude total probability $\delta>0$, whence
$B_3\leq1-\delta$ and $e^{h_w}\leq(1-\delta)^{1/3}<1$.
Set $q_w=e^{h_w}$ and use the cyclic sandwich.
The fixed matchings used in Corollary~\ref{cor:prefix} have positive
probability under full support, giving the two smaller sizes.
\end{proof}

For example, support on rankings with offset $+1$ first and $-1$ second
gives every even cycle a stable adjacent-pair matching, so $p_{2j}=1$:
each matched vertex has its first or second choice, and every vertex
with its second choice rejects both distance-two neighbors.
Full support cannot simply be removed from Corollary~\ref{cor:weighted}.
This is an obstruction to that unposed generalization, not a
counterexample to Mertens's uniform-law conjecture.

\begin{corollary}[Triangle packing]\label{cor:packing}
Let $G$ be a finite simple graph on $N$ vertices with maximum degree at
most $D\geq2$, and let $t$ be the number of triangles in any
vertex-disjoint triangle family. Under independent uniform neighbor
rankings, its stable-partial-matching probability $P(G)$ satisfies
\[
 D^{-N}\leq P(G)\leq
 \left(1-\frac{2}{[D(D-1)]^3}\right)^t.
\]
In particular $p_N\leq(1-1/864)^{\lfloor N/3\rfloor}$ for range-two
cycles. More generally an arbitrary vertex-disjoint family $\mathcal C$
of odd cycles gives the upper bound
\[
 P(G)\leq\prod_{C\in\mathcal C}
 \left(1-\frac{2}{[D(D-1)]^{|C|}}\right).
\]
These are two exponential bounds, not a common-base conclusion
for arbitrary graph families.
\end{corollary}
\begin{proof}
Choose a maximal matching and put its assigned partner first at every
matched vertex. Its unmatched vertices are independent, so every such
profile is stable. The event has probability
$\prod_{v\text{ matched}}\deg(v)^{-1}\geq D^{-N}$.
For a triangle with vertices $u,v,z$, the two exclusive orientations
have total probability
$2/\prod_{x\in\{u,v,z\}}[\deg(x)(\deg(x)-1)]$,
at least $2/[D(D-1)]^3$, and each prevents stability.
The events on vertex-disjoint triangles are independent.
For a range-two cycle, take the $\lfloor N/3\rfloor$ consecutive triples.
For any odd cycle the same exclusive-order construction prevents a
stable matching: an internally unmatched vertex and its successor form
a blocking pair. Its two orientations have probability
$2/\prod_{v\in C}[\deg(v)(\deg(v)-1)]$. Independence for disjoint
vertex sets gives the product bound.
\end{proof}

\begin{corollary}[Mertens's range-one formula]\label{cor:rangeone}
For the range-one cycle on $N\geq3$ vertices with uniform independent
rankings, the solvability probability is $1$ for even $N$ and
$1-2^{1-N}$ for odd $N$, as in \cite[equation~(13)]{mertens2005}.
\end{corollary}
\begin{proof}
If the first choices are not all directed around the cycle in the same
orientation, some adjacent pair ranks each other first. Fix that pair;
the remaining graph is a path. On a finite forest with an edge, directing
each nonisolated vertex toward its first choice produces a mutual-first
pair, since a directed cycle of length at least three would yield an
undirected cycle. Repeatedly fix and delete such pairs. Earlier pairs
reject every later vertex, so the resulting matching is stable.
If all choices have one orientation, an odd cycle is exclusive and
unsolvable; an even cycle has a stable alternating adjacent-pair matching.
The two exclusive full-cycle profiles each have probability $2^{-N}$.
\end{proof}
Thus the range-one exception precedes Conjecture~3 in its source;
applying exponential decay with $q<1$ to that exception would not be the
claim posed there.

\begin{corollary}[Complete-graph endpoint; Young--Studebaker]
\label{cor:complete}
At $p=1$, the even-$N$ solvability probability in Mertens's
Conjecture~5, equation~(19), is
\[
 P_N\sim C_*N^{-1/4},\qquad
 C_*=\frac{e\,2^{1/4}\Gamma(3/4)}{\sqrt\pi},
\]
by \cite[Theorem~1.1]{young2026}. Here
$\Gamma(s)=\int_0^\infty t^{s-1}e^{-t}\,dt$ for $s>0$.
The predicted coefficient $e\sqrt{2/\pi}$ in equations~(7) and~(19)
is corrected by this theorem; the interior $0<p<1$ is not covered.
\end{corollary}
\begin{proof}
The graph $G(N,1)$ is deterministically the complete graph.
For even $N$ a stable partial matching in a complete graph cannot have
singles: two singles would block, and their number must be even.
The models therefore coincide, and the cited theorem applies.
\end{proof}

\section{Coverage and verification}
The written proofs above establish the all-length statement and the
finite-prefix extension. The following pinned Lean declarations give
formal counterparts with the same labelled-profile and single-agent
conventions; their namespace prefix is
\texttt{D5.S3.Combinatorics.Graph.Roommates}.

\begin{center}\small
\begin{tabular}{@{}p{0.44\linewidth}p{0.50\linewidth}@{}}
\toprule
Declaration & Mathematical coverage \\
\midrule
\leanfile{SourceModel}{\texttt{SourceModel.card\_profile}} &
Independent uniform rankings; $24^{n+5}$ profiles.\\[3pt]
\leanfile{OpenCertificates}{\texttt{OpenCertificates.count\_lower},\newline
\texttt{count\_submultiplicative},\newline\texttt{count\_three\_bound}} &
$6^m\leq b_m$, submultiplicativity, and $b_3\leq13808$.\\[3pt]
\leanfile{BufferLayout}{\texttt{BufferLayout.buffer\_completion}} &
Six-site free-vertex matching with at most one central single.\\[3pt]
\leanfile{PaddedCompletion}{\texttt{PaddedCompletion.}\newline
\texttt{count\_padded\_supermultiplicative}} &
$b_mb_k\leq b_{m+k+6}$.\\[3pt]
\leanfile{CyclicClosure}{\texttt{CyclicClosure.}\newline
\texttt{count\_le\_successful\_profile\_count},\newline
\texttt{source\_count\_le\_open\_count}} &
$b_m\leq A_{m+6}$ and $A_N\leq b_N$.\\[3pt]
\leanfile{GrowthTransfer}{\texttt{GrowthTransfer.result}} &
$\exists q,c,C>0$, $q<1$, $cq^{n+5}\leq p_{n+5}\leq Cq^{n+5}$
for every $n\geq2$, with no combinatorial premises.\\
\bottomrule
\end{tabular}
\end{center}

The accompanying exact-arithmetic checker covers open words of lengths
zero through three, their cyclic completions at sizes seven through nine,
every gap-occupancy pattern, selected fully supported rational laws, and
range-one cycles of sizes three through nine. Its detailed counts and
scope are in the supporting files. The all-length injections, base
enclosures, weighted extension, triangle-packing bounds and range-one
formula are written proofs. The complete-graph asymptotic is the cited
theorem of Young and Studebaker, not a consequence of the range-two
argument. The pinned declarations above concern the uniform range-two
model at $N\geq7$; the extension to $N=5,6$, explicit interval for $q$,
and Corollaries~\ref{cor:enclosures}--\ref{cor:complete} have the written
or cited coverage just specified.

For context, the following table separates the statements in the posing
paper \cite{mertens2005}; ``not covered'' describes the scope of this
note rather than a claim about their current status.
\begin{center}\small
\begin{tabular}{@{}p{0.30\linewidth}p{0.64\linewidth}@{}}
\toprule
Source item & Coverage \\
\midrule
Conjecture 1, equation (7) & Prior proof cited: Young--Studebaker (2026)
prove the exponent and correct the asserted coefficient.\\[3pt]
Conjecture 2, equation (9) & Not covered: expected size of the odd stable
cycles in unsolvable complete-graph instances.\\[3pt]
Conjecture 3, equation (15) & Proved here for $d=1,r=2$.
Other dimensions and ranges are not covered.\\[3pt]
Equation (13) & Reproved; prior proof Mertens (2005).
Corollary~\ref{cor:rangeone}; probability tends to one.\\[3pt]
Five-agent range-two cycle & Prior exact value Mertens (2015),
\cite[equation~(21b)]{mertenssmall2015}: $p_5=4075/6912$.\\[3pt]
Equation (14) and following heuristic & Proved: independent triangle
and odd-cycle packing gives Corollary~\ref{cor:packing}, without asserting a common
base for other grids.\\[3pt]
Conjecture 4, equations (17)--(18) & Not covered: sparse
Erd\H{o}s--R\'enyi acceptability graphs.\\[3pt]
Conjecture 5, equation (19) & Refuted-and-corrected at $p=1$ by the
cited Young--Studebaker coefficient; $0<p<1$ not covered.\\[3pt]
Weighted rankings; effective base bounds & Not posed; proved in
Corollaries~\ref{cor:weighted} and~\ref{cor:enclosures}.
No full-support claim is made for arbitrary restricted laws.\\
\bottomrule
\end{tabular}
\end{center}

\begin{thebibliography}{9}
\bibitem{mertens2005}
S. Mertens, Random stable matchings,
\emph{Journal of Statistical Mechanics: Theory and Experiment} (2005),
P10008. \href{https://doi.org/10.1088/1742-5468/2005/10/P10008}{doi:10.1088/1742-5468/2005/10/P10008}.
Posing locators refer to \href{https://arxiv.org/abs/cond-mat/0509221v2}{arXiv:cond-mat/0509221v2}.
\bibitem{mertens2015}
S. Mertens, Stable roommates problem with random preferences,
\emph{Journal of Statistical Mechanics: Theory and Experiment} (2015),
P01020. \href{https://doi.org/10.1088/1742-5468/2015/01/P01020}{doi:10.1088/1742-5468/2015/01/P01020};
\href{https://arxiv.org/abs/1401.5269v2}{arXiv:1401.5269v2}.
\bibitem{mertenssmall2015}
S. Mertens, Small random instances of the stable roommates problem,
\emph{Journal of Statistical Mechanics: Theory and Experiment} (2015),
P06034. \href{https://doi.org/10.1088/1742-5468/2015/06/P06034}{doi:10.1088/1742-5468/2015/06/P06034};
\href{https://arxiv.org/abs/1502.06635v1}{arXiv:1502.06635v1},
Section~5, equation~(21b).
\bibitem{young2026}
C. Young and I. Studebaker, Sharp asymptotics for the solvability
probability of random stable roommates,
\href{https://arxiv.org/abs/2610.07741v1}{arXiv:2610.07741v1} (2026), Theorem~1.1.
\bibitem{chin2026}
B. Chin and M. Michelen, The random stable roommates problem typically
has no solution, \href{https://arxiv.org/abs/2601.07612v1}{arXiv:2601.07612v1}
(2026), Theorem~1.
\end{thebibliography}
\smallskip
\noindent\small\copyright\ 2026 Haobo Ma and Wenlin Zhang.
Licensed under Creative Commons Attribution 4.0 International
(\href{https://creativecommons.org/licenses/by/4.0/}{CC BY 4.0}).
\end{document}
